\documentclass[11pt,a4paper]{article}
\usepackage[utf8]{inputenc}
\usepackage[T1]{fontenc}
\usepackage{lmodern}
\usepackage{amsmath}
\usepackage[margin=2cm]{geometry}
\usepackage{enumitem}
\pagestyle{plain}
\setlength{\parindent}{0pt}
\setlength{\parskip}{0pt}
\begin{document}
{\Large\bfseries Tarea de oxiácidos y oxisales}
\vspace{12pt}
\begin{enumerate}[label=\textbf{\arabic*.},leftmargin=*,itemsep=11pt,topsep=0pt]
\item \begin{minipage}[t]{\linewidth}
¿Cuál es la fórmula del ácido perclórico?

\vspace{5pt}
\begin{tabular}{@{}p{0.47\linewidth}p{0.47\linewidth}@{}}
a)~$\mathrm{HClO_{4}}$ & b)~$\mathrm{HClO_{2}}$ \\
c)~HClO & d)~$\mathrm{HClO_{3}}$ \\
\end{tabular}
\end{minipage}

\item \begin{minipage}[t]{\linewidth}
¿Cuál es la fórmula del ácido sulfuroso?

\vspace{5pt}
\begin{tabular}{@{}p{0.47\linewidth}p{0.47\linewidth}@{}}
a)~$\mathrm{HSO_{3}}$ & b)~$\mathrm{H_{2}SO_{2}}$ \\
c)~$\mathrm{H_{2}SO_{3}}$ & d)~$\mathrm{H_{2}SO_{4}}$ \\
\end{tabular}
\end{minipage}

\item \begin{minipage}[t]{\linewidth}
Si el fósforo con +5 forma el ácido con tres moléculas de agua, ¿cuál es la fórmula del ácido fosfórico?

\vspace{5pt}
\begin{tabular}{@{}p{0.47\linewidth}p{0.47\linewidth}@{}}
a)~$\mathrm{HPO_{4}}$ & b)~$\mathrm{H_{3}PO_{3}}$ \\
c)~$\mathrm{H_{3}PO_{4}}$ & d)~$\mathrm{H_{2}PO_{4}}$ \\
\end{tabular}
\end{minipage}

\item \begin{minipage}[t]{\linewidth}
¿Cuál es el nombre tradicional del $\mathrm{HNO_{2}}$?

\vspace{5pt}
\begin{tabular}{@{}p{0.47\linewidth}p{0.47\linewidth}@{}}
a)~ácido nitroso & b)~ácido pernítrico \\
c)~ácido nítrico & d)~ácido hiponitroso \\
\end{tabular}
\end{minipage}

\item \begin{minipage}[t]{\linewidth}
¿Cuál es el nombre sistemático del $\mathrm{H_{2}CO_{3}}$?

\vspace{5pt}
\begin{tabular}{@{}p{0.47\linewidth}p{0.47\linewidth}@{}}
a)~trioxocarbonato (IV) de hidrógeno & b)~dioxocarbonato (IV) de hidrógeno \\
c)~ácido trioxocarbónico (IV) & d)~carbonato de hidrógeno (II) \\
\end{tabular}
\end{minipage}

\item \begin{minipage}[t]{\linewidth}
¿Cuál es el nombre stock del $\mathrm{HBrO_{3}}$?

\vspace{5pt}
\begin{tabular}{@{}p{0.47\linewidth}p{0.47\linewidth}@{}}
a)~ácido dioxobrómico (III) & b)~trioxobromato (V) de hidrógeno \\
c)~ácido trioxobrómico (V) & d)~ácido trioxobromoso (V) \\
\end{tabular}
\end{minipage}

\item \begin{minipage}[t]{\linewidth}
¿Cuál es la fórmula del hipoclorito de sodio?

\vspace{5pt}
\begin{tabular}{@{}p{0.47\linewidth}p{0.47\linewidth}@{}}
a)~NaCl & b)~NaClO \\
c)~$\mathrm{NaClO_{3}}$ & d)~$\mathrm{NaClO_{2}}$ \\
\end{tabular}
\end{minipage}

\item \begin{minipage}[t]{\linewidth}
¿Cuál es la fórmula del carbonato de calcio?

\vspace{5pt}
\begin{tabular}{@{}p{0.47\linewidth}p{0.47\linewidth}@{}}
a)~$\mathrm{Ca(CO_{3})_{2}}$ & b)~$\mathrm{CaCO_{3}}$ \\
c)~$\mathrm{CaCO_{2}}$ & d)~$\mathrm{Ca_{2}CO_{3}}$ \\
\end{tabular}
\end{minipage}

\item \begin{minipage}[t]{\linewidth}
¿Cuál es la fórmula del permanganato de potasio?

\vspace{5pt}
\begin{tabular}{@{}p{0.47\linewidth}p{0.47\linewidth}@{}}
a)~$\mathrm{K(MnO_{4})_{2}}$ & b)~$\mathrm{KMnO_{3}}$ \\
c)~$\mathrm{K_{2}MnO_{4}}$ & d)~$\mathrm{KMnO_{4}}$ \\
\end{tabular}
\end{minipage}

\item \begin{minipage}[t]{\linewidth}
¿Cuál es el nombre sistemático del $\mathrm{Na_{2}SO_{3}}$?

\vspace{5pt}
\begin{tabular}{@{}p{0.47\linewidth}p{0.47\linewidth}@{}}
a)~bis[trioxosulfato (IV)] de sodio & b)~tetraoxosulfato (VI) de disodio \\
c)~trioxosulfato (IV) de sodio & d)~trioxosulfato (IV) de disodio \\
\end{tabular}
\end{minipage}

\end{enumerate}
\end{document}
